% TOP_PROBLEM: {"id": 3198, "title": "Acyclic edge-colouring", "category_id": 3, "category": {"id": 3, "name": "graph_theory", "display_name": "Graph Theory", "description": "Problems involving graphs, networks, and their properties.", "slug": "graph-theory", "order_index": 3, "created_at": "2026-07-31T15:26:25.671Z"}, "status": "open", "problem_number": "OPG-145", "set_id": 12}
% FIRST_POSTED: 2026-10-01
% ENTRY_KIND: statement_only
% UnsolvedMath ID 3198; source catalogue code OPG-145
% !TeX program = lualatex
% Definition and sources only; this file is not an LLM solution attempt.
\documentclass[11pt,a4paper]{article}
\usepackage[margin=25mm]{geometry}
\usepackage{fontspec}
\setmainfont{lmroman10-regular.otf}[BoldFont=lmroman10-bold.otf,
 ItalicFont=lmroman10-italic.otf,BoldItalicFont=lmroman10-bolditalic.otf]
\usepackage{amsmath,amssymb,mathtools}
\usepackage{unicode-math}
\setmathfont{latinmodern-math.otf}
\AtBeginDocument{\let\setminus\smallsetminus}
\usepackage{xurl,hyperref}
\hypersetup{colorlinks=true,urlcolor=blue}
\setcounter{secnumdepth}{0}
\setlength{\parindent}{0pt}
\setlength{\parskip}{5pt}
\setlength{\emergencystretch}{3em}
\begin{document}
\section*{Acyclic edge-colouring}\label{3198}
{\small UnsolvedMath ID 3198 · OPG-145 · Graph Theory · OpenGarden}
\subsection{Definitions and mathematical statement}
Let $G=(V,E)$ be a finite simple undirected graph,
and put $\Delta=\max_{v\in V}\deg_G(v)$, with
$\Delta=0$ for the empty graph. A proper
$q$-edge-colouring is a function
$c:E\to\{1,\ldots,q\}$ assigning different
colours to any two edges sharing an endpoint.
Such a colouring is acyclic if every simple
cycle in $G$ uses at least three colours.

Equivalently, for every distinct pair of
colours $a,b$, the spanning subgraph with
edge set $c^{-1}(\{a,b\})$ is a forest.
Properness forces a two-coloured cycle,
if present, to alternate its colours and
therefore have even length. Eliminating
all such cycles is the additional acyclic
requirement.

The conjecture states that every finite
simple graph has an acyclic proper edge
colouring using at most $\Delta+2$ colours:
\[
 \exists c:E\to\{1,\ldots,\Delta+2\}
 \quad c\text{ proper and every cycle uses at least three colours}.
\]
Unused colours are allowed. The bound uses
the maximum degree of the original graph,
and is independent of its order, girth,
or connectedness.

The word acyclic refers to the union of
each pair of colour classes, not to the
whole graph, which may contain many cycles.
A proper edge colouring already makes
each individual colour class a matching;
the difficulty is controlling their pairwise
unions simultaneously with only two colours
beyond the degree bound.
\subsection{Short English statement}
Can every simple graph be properly edge-coloured with at most maximum degree plus two colours so that no cycle uses only two colours?
\subsection{Sources}
\begin{itemize}
\item[S1] ULAM.AI and the UnsolvedMath Contributors, supplied problems.json, record 3198 (OPG-145), OpenGarden collection. Catalogue identity and source transcription; CC BY 4.0.
\url{https://huggingface.co/datasets/ulamai/UnsolvedMath}.
\item[S2] Open Problem Garden, Acyclic edge-colouring. Original problem and discussion; the mathematical formulation below is expanded from the supplied source record.
\url{http://www.openproblemgarden.org/op/acyclic_edge_coloring}.
\end{itemize}
\end{document}
